\documentclass[12pt]{article}

\usepackage{amsmath, amssymb, amsthm}
\usepackage{geometry}
\usepackage{graphicx}
\usepackage{booktabs}
\usepackage{tikz}
\usepackage{hyperref}
\usepackage{float}
\usepackage{listings}
\usepackage{xcolor}
\usepackage{tikz}
\usetikzlibrary{arrows.meta}
\numberwithin{equation}{section}
\usepackage{hyperref}

\title{One Pricer a Day}
\author{Neel Dev Sen}
\date{July 2026}

\lstdefinestyle{cpp}{
	language=C++,
	basicstyle=\ttfamily\small,
	keywordstyle=\color{blue},
	commentstyle=\color{gray},
	stringstyle=\color{purple},
	numbers=left,
	numberstyle=\tiny\color{gray},
	stepnumber=1,
	frame=single,
	breaklines=true,
	showstringspaces=false
	tabsize=4
}

\lstdefinestyle{python}{
	language=Python,
	basicstyle=\ttfamily\small,
	keywordstyle=\color{blue},
	commentstyle=\color{gray},
	stringstyle=\color{purple},
	numbers=left,
	numberstyle=\tiny\color{gray},
	stepnumber=1,
	frame=single,
	breaklines=true,
	showstringspaces=false
}

\begin{document}
	
	\maketitle
	\newpage
	\tableofcontents
	\newpage

\section{Introduction}
This is a project where I develop derivative pricers using \textbf{C++}. My goal is to do at least one a day and you can access all of my \href{https://github.com/neeldevsen/one-pricer-a-day}{Github Repo}
\newpage
\section{European Options}
\subsection{Day 1: Black Scholes}
The first option that I am going to be pricing is the European option using the closed-form Black-Scholes equation. \newline \newline In \textbf{Github} the full \textbf{C++} script is on this link:
\href{https://github.com/neeldevsen/one-pricer-a-day/blob/main/day1-BlackScholes.cpp}{Day1-BlackScholesPricer}
\subsubsection{The Black Scholes PDE}
  The Black Scholes equation is a partial differential equation that is used to model how option prices change as the initial stock price, time to maturity, risk-free interest rate and implied volatility change. The partial differential equation is shown below:
  \begin{equation} 
  	\frac{\partial V}{\partial t} + \frac{1}{2} \sigma^2 S^2 \frac{\partial^2 V}{\partial S^2} + rS\frac{\partial V}{\partial S} - rV = 0
 \end{equation} \newline
 In this equation $V$ is the option's price, $\sigma$ is the implied volatility, $r$ is the risk-free rate of interest and $t$ is the time to maturity. \newline \newline $\frac{\partial V}{\partial t} $ is the rate of change of the option's price to time, $\frac{\partial V}{\partial S}$ is the rate of change of the option's price to the underlying stock price and  $\frac{\partial^2 V}{\partial S^2}$ is a concavity factor on how the option price changes with the stock price.  When solving the Black-Scholes there are two cases: when the option is a call option, or when the option is a put option.
 \newpage
 \subsubsection{Closed-Form Solution for European Call Options}
 The closed-form solution for the Black-Scholes equation where the option is a call option is:
 \begin{equation}
 	C(t,S) = S\Phi(d_1) - Ke^{-r(T-t)}\Phi(d_2)
 	\label{eq:2.2}
 \end{equation}
 where: 
 \begin{equation*}
 	d_1 = \frac{\ln(\frac{S}{K}) + (r + \frac{1}{2}\sigma^2)(T-t)}{\sigma\sqrt{T-t}}
 \end{equation*}
  and:
  \begin{equation*}
  	d_2 = d_1 - \sigma\sqrt{T-t}
  \end{equation*}
  In this equation $T-t$ represents the time to maturity and $K$ represents the strike price. \newline \newline
  $\Phi(z)$ is the standard normal cumulative distribution function which can be represented with the error function $\operatorname{erf}(z)$: 
  \begin{equation}
  	\Phi(z) = \frac{1}{2}(1 + \operatorname{erf}(\frac{z}{\sqrt{2}}))
  	\label{eq:2.3}
  \end{equation}
 \newline In C++ this can be implemented as:
  
 \begin{lstlisting}[style=cpp]
#include <iostream>
#include <cmath>
#include <type_traits>
 
template <typename T>
auto normalCDF(T x) -> std::common_type_t <double , T>
{
	using commonType = std::common_type_t<double, T>;
	return static_cast<commonType>(0.5) * (static_cast<commonType>(1) + std::erf(x / std::sqrt(2.0)));
}
 \end{lstlisting}
 The function we are creating is called \textbf{normalCDF}.This listing uses the \textbf{cmath} header file in order to import \textbf{std::erf}. It also uses a function template with a type placeholder \textbf{T} so that we can use this function for any type given. In order to ensure that this function has at least a type \textbf{double} we use the header file \textbf{type traits} which allows us to specify that we want this function to be at least double in line 6 trailing return type. In line 8 we ensure that the \textbf{commonType} alias used amongst all the literals and variables in this function are also at least of type \textbf{double}. This function returns the result shown in equation \ref{eq:2.3} with a type of at least \textbf{double}. 
 \newpage
 Now to implement the function shown in equation \ref{eq:2.2} in C++ we use: 
 \begin{lstlisting}[style=cpp]
 template <typename SType, typename KType, typename rType, typename sigmaType, typename tType>
 auto blackScholesCall(SType S, KType K, rType r, sigmaType sigma, tType t) -> std::common_type_t<double, SType, KType, rType, sigmaType, tType>
 {
 	using commonType = std::common_type_t<double, SType, KType, rType, sigmaType, tType>;
 	auto S_new {static_cast<commonType>(S)};
 	auto K_new {static_cast<commonType>(K)};
 	auto r_new {static_cast<commonType>(r)};
 	auto sigma_new {static_cast<commonType>(sigma)};
 	auto t_new {static_cast<commonType>(t)};
 		
 		
 	auto d1 {(std::log(S_new / K_new) + (r_new + 0.5 * sigma_new * sigma_new) * t_new) / (sigma_new * std::sqrt(t_new))};
 	auto d2 {d1 - sigma_new * std::sqrt(t_new)};
 		
 		
 	auto C {S_new * normalCDF(d1) - K_new * std::exp(-r_new * t_new) * normalCDF(d2)};
 	return C;
 	}
 \end{lstlisting}
 The function in this listing is called \textbf{blackScholesCall}. This listing again uses the same header files and the same principles as the one before and converts every thing into a \textbf{commonType} of at least type \textbf{double}. The formula for $d_1$ is shown in line 12 and the formula for $d_2$ is shown in line 13.  Finally the result in equation \ref{eq:2.2} is shown in line 16 using the previously defined function \textbf{normalCDF}. The function returns the value of the \textbf{call option} with at least type \textbf{double}. 
 \newpage
\subsubsection{Closed-Form Solution for European Put Options}
 The closed-form solution for the Black-Scholes equation where the option is a put option is:
\begin{equation}
	P(t,S) = Ke^{-r(T-t)}\Phi(-d_2) - S\Phi(-d_1)
	\label{eq:2.4}
\end{equation}
using the same definitions of $d_1$ and $d_2$ in \ref{eq:2.2} \newline
The implementation of this in C++ is:

 \begin{lstlisting}[style=cpp]
template <typename SType, typename KType, typename rType, typename sigmaType, typename tType>
auto blackScholesPut(SType S, KType K, rType r, sigmaType sigma, tType t) -> std::common_type_t<double, SType, KType, rType, sigmaType, tType>
{
	using commonType = std::common_type_t<double, SType, KType, rType, sigmaType, tType>;
	auto S_new {static_cast<commonType>(S)};
	auto K_new {static_cast<commonType>(K)};
	auto r_new {static_cast<commonType>(r)};
	auto sigma_new {static_cast<commonType>(sigma)};
	auto t_new {static_cast<commonType>(t)};
	
	
	auto d1 {(std::log(S_new / K_new) + (r_new + 0.5 * sigma_new * sigma_new) * t_new) / (sigma_new * std::sqrt(t_new))};
	auto d2 {d1 - sigma_new * std::sqrt(t_new)};
	
	
	auto P {K_new * std::exp(-r_new * t_new) * normalCDF(-d2) - S_new * normalCDF(-d1)};
	return P;
}
\end{lstlisting}
 The function in this listing is called \textbf{blackScholesPut}. This listing is almost identical to the one for \textbf{blackScholesCall}. The formula for $d_1$ is again shown in line 12 and the formula for $d_2$ is shown in line 13.  Finally the result in equation \ref{eq:2.4} is shown in line 16 using the function \textbf{normalCDF}. The function returns the value of the \textbf{put option} with at least type \textbf{double}. 
\newpage
\end{document}
